\CQUsection{总结与展望}

\CQUsubsection{研究工作总结}
本文围绕颅内动脉瘤智能分析问题，构建了由病灶分割与破裂风险预测组成的研究框架，并在统一实验流程下完成了方法设计、模型训练、结果评估与原型验证。研究工作的核心目标在于：一方面，利用三维医学影像分割方法实现对动脉瘤区域的精细识别与结构刻画；另一方面，结合病例层面的临床信息、影像表型信息及分割衍生结构特征，建立面向破裂风险评估的分类模型，从而形成由解剖结构分析到病例级风险判别的完整技术链条。

在方法层面，本文针对 CTA 图像中动脉瘤体积小、背景复杂和类别不平衡等问题，构建了以三维 Attention U-Net 为主干的分割模型，并结合多尺度上下文增强与方向敏感注意力机制提升病灶表征能力。针对风险预测任务中多源异构表格特征之间存在复杂非线性交互的问题，本文进一步设计了基于字段嵌入与交叉注意力的风险模型，以增强对变量间高阶依赖关系的刻画能力。通过上述设计，本文在方法上实现了影像空间表征与病例级决策建模的有机衔接。

\CQUsubsection{主要研究结论}
基于实验结果，可以得到以下几点主要结论。

第一，在动脉瘤分割任务中，所构建的三维分割模型在验证阶段表现出较高的体素级判别能力和较稳定的阳性检出能力，说明该模型能够较好地从复杂 CTA 背景中识别目标病灶区域。注意力门控机制有助于抑制无关背景响应，多尺度上下文聚合与方向敏感注意力则进一步增强了模型对小目标、复杂形态和空间走向特征的表达能力。

第二，在破裂风险预测任务中，基于交叉注意力的深度表格模型能够有效学习临床变量、影像表型变量与结构化形态特征之间的交互关系。测试结果表明，该模型在 F1 与 AUC 等核心指标上接近最优融合模型，并整体优于多种传统单模型基线，说明通过字段级表示学习和注意力交互机制能够提升多源医学表格数据的风险判别能力。

第三，融合模型取得了最高的综合性能，这表明深度注意力模型与传统集成学习模型在当前任务中具有一定互补性。前者更擅长挖掘连续隐空间中的复杂关联模式，后者则在中小样本场景下表现出较强稳健性。因此，在医学风险预测问题中，单一模型范式未必能够完全占优，合理结合不同方法的优势可能是提升性能的重要方向。

第四，从整体流程角度看，本文实现了由病灶分割到风险评估的闭环验证，说明基于影像结构识别结果构建病例级风险表征具有可行性。这一结果表明，分割任务不仅可以作为独立的病灶检测工具，也可以为后续风险判别提供具有解剖学意义的结构先验，从而增强模型的医学解释基础。

\CQUsubsection{研究不足}
尽管本文在方法构建与实验验证方面取得了一定结果，但仍存在若干不足，主要体现在以下几个方面。

首先，数据规模与来源范围仍然有限。医学影像与临床数据的采集和标注成本较高，导致当前研究更多是在相对有限的数据条件下开展。样本量不足不仅会限制深度模型对复杂模式的充分学习，也可能使实验结果对特定数据划分较为敏感，从而影响统计稳健性与泛化能力。

其次，外部验证仍然不足。本文实验主要基于既定数据划分展开，虽然能够反映模型在当前数据分布下的有效性，但尚不能充分说明模型在不同中心、不同扫描设备、不同人群结构和不同标注标准下的推广能力。对于医学人工智能方法而言，外部独立验证是评价模型真实应用价值的重要环节。

再次，风险预测任务仍然面临类别不平衡与特征异质性带来的挑战。虽然本文通过类别加权、阈值搜索和注意力交互机制在一定程度上缓解了相关问题，但在小样本条件下，深度表格模型对特征噪声、超参数设置和训练随机性的敏感性仍然不可忽视。这也是为何传统树模型在当前实验中依然具有较强竞争力的原因之一。

最后，本文的结果解释仍主要集中于性能指标层面，对模型内部决策依据的临床可解释分析仍有待加强。例如，不同临床变量、形态变量和影像表型变量在风险预测中的相对贡献尚需通过更系统的可解释性方法进行量化；分割结果中的局部误差如何进一步影响后续风险建模，也值得深入研究。

\CQUsubsection{未来工作展望}
针对上述不足，未来研究可从以下几个方向进一步推进。

第一，扩大数据规模并开展多中心验证。后续工作可进一步引入来自不同中心、不同设备和不同扫描协议的数据，以增强样本多样性，并通过外部测试集、交叉验证和置信区间分析更严格地评估模型稳定性与泛化性能。这不仅有助于提升方法的可信度，也更符合医学人工智能向临床应用转化的基本要求。

第二，进一步优化分割模型的结构表达与监督机制。未来可探索更强的三维上下文建模策略、边界敏感损失函数、层级监督机制以及面向小目标的困难样本挖掘方法，以提升病灶边界恢复能力和对复杂解剖背景的鲁棒性。同时，还可研究如何将血管拓扑先验或形态约束显式引入模型训练，以增强分割结果的结构合理性。

第三，深化风险预测中的特征融合与可解释性建模。后续可结合更丰富的形态学定量指标、时序随访信息或多模态影像表征，构建更全面的病例级特征空间；同时可引入注意力可视化、特征归因、概率校准和不确定性估计等方法，从而提高风险预测结果的可解释性、可信度与临床可读性。

第四，探索分割与风险预测的联合优化机制。本文采用的是两阶段分析框架，具有结构清晰、实现简洁的优点，但分割与风险建模之间仍可能存在更深层次的协同关系。未来可以考虑建立端到端联合学习框架，使影像空间特征、病灶结构表征与风险标签在统一目标下协同优化，从而进一步增强任务间的信息传递效率。

第五，推动方法从实验验证走向更规范的应用研究。未来工作可进一步完善推理效率评估、结果一致性分析、临床工作流适配和人机协同机制设计，并结合前瞻性研究验证模型在真实场景中的辅助价值。只有在算法性能、可解释性、稳定性和使用流程等多个层面均得到充分论证，相关研究成果才可能具备更高的临床转化潜力。

\CQUsubsection{结束语}
总体而言，本文围绕颅内动脉瘤智能分析任务，完成了从三维病灶分割到破裂风险预测的系统研究，并在方法设计与实验分析层面取得了较为完整的阶段性结果。本文工作表明，结合多尺度医学影像表征与多源临床特征融合，可以为动脉瘤的自动识别与风险评估提供具有一定实证基础的技术路径。尽管当前研究仍存在数据规模、外部验证和可解释性等方面的局限，但其在方法框架与实验结果上已为后续深入研究奠定了基础，也为相关医学人工智能技术的持续优化与应用探索提供了参考。